\section{外推边界：其他机制、LSTM 与 Scaling Laws}

Induction-head story 很有吸引力，因为它把行为、数学结构与 intervention 连在一起；越是整洁的故事，越需要主动寻找不能解释的部分。讲者在结尾明确允许其他 mechanisms，并把许多大模型计算留作开放问题。

\subsection{In-context nearest neighbor 不是全部 meta-learning}

Olah 将 induction 压缩为 in-context nearest neighbor，同时指出 soft heads 可以学习 functions。这个描述适合解释从 examples 中匹配和延续 pattern，却不必然覆盖 gradient-descent-like inner optimization、Bayesian updating、multi-step reasoning 或工具使用。

\lecturefigure{slide-61-in-context-nearest-neighbor.jpg}{Induction 作为 in-context nearest neighbor；soft heads 可以学习函数关系。}{Stanford Online 官方视频 00:55:30--00:55:33。}

\lecturefigure{slide-62-alternative-contributors.jpg}{其他 mechanisms 可能同时贡献 large-model ICL。}{Stanford Online 官方视频 00:55:34--00:55:59。}

\begin{warningbox}{漂亮的 mechanistic story 仍可能不完备}
Induction heads 可以解释某类 lookup/copy behavior，并在 small models 中解释 chosen ICL score；大型模型可能同时使用 MLP memory、distributed feature binding、other heads、implicit optimization 或更长 path。最好的结论是“一个重要机制”，不是“唯一机制”。
\end{warningbox}

\teachervoice{课堂结尾给出的 present theory 大致包含两类算法：一类是广泛的“看到过某词，未来合适位置可能再出现”的 copying；另一类是 induction-head nearest neighbor。讲者保留了 “possibly other things” 的空间。}

\subsection{为什么 Transformer 比 bounded-state LSTM 更适合长 context lookup}

Attention 可以直接访问任意过去 position，检索成本主要由 context length 与 attention implementation 决定；传统 LSTM 必须把所有历史压进固定维 recurrent state，长距离精确 lookup 容易衰减。这个差异使 induction-style retrieval 在 Transformer 中更自然，但不意味着 LSTM 数学上无法表达任何 meta-learning，也不意味着 attention 自动学到正确 lookup。

\lecturefigure{slide-63-transformer-vs-lstm.jpg}{长 context ICL：attention 的直接检索路径与 LSTM 固定状态瓶颈。}{Stanford Online 官方视频 00:56:00--00:56:21。}

\begin{knowledgebox}{不要把架构比较说成绝对不可能}
有限精度、有限 hidden state 和实际训练条件下，LSTM 长距离保真 retrieval 很困难；理论上更大状态、外部 memory 或特殊训练可改变结论。课堂中的 “impossible” 应按工程语境理解为该类 mechanism 不自然、难以稳定扩展。
\end{knowledgebox}

\subsection{Induction heads 会解释 meta-learning scaling laws 吗}

若 ICL 近似 context nearest neighbor，那么 capability 可能随 training data 中 pattern diversity、representation quality 与 context retrieval 共同变化，而不只由 parameter count 决定。讲者提出这可能帮助理解 scaling-law 曲线，但没有给出完整理论或验证。

\lecturefigure{slide-64-meta-learning-scaling-laws.jpg}{开放问题：induction-head mechanism 是否能解释 meta-learning scaling laws？}{Stanford Online 官方视频 00:56:22--00:56:47。}

\teachervoice{Olah 对 scaling-law 解释明显更 tentative：这只是可能的路线。讲义不应把一张问题页写成已建立的 law，更不能用 induction story 取代对 data distribution、optimization 与 architecture scaling 的实证分析。}

\subsection{真正尚未理解的部分}

Q\&A 中 Olah 列出 MLP layers、arithmetic、multi-speaker/dialogue representation 与大型模型的大量内部机制。Attention-only toy models 的 near-complete account 是方法演示，不是研究终点；大模型的 nonlinear feature superposition、normalization 与 distributed circuits 仍然困难。

\teachervoice{“还有太多不了解，以至于很难回答下一步是什么”是本讲最诚实的结尾。Mechanistic interpretability 的成熟度应由可复现 circuit、干预结果与失败案例衡量，而不是由术语数量衡量。}

\subsection{本章小结}

Induction heads 是有结构、有 intervention 支持的重要 ICL mechanism，但大型模型可能使用其他 algorithms。Transformer 为长 context retrieval 提供自然接口，scaling-law 连接仍属开放假说，而 MLP 与 distributed circuits 是下一阶段的核心困难。

